(1) 初态总能量 \(E+mc^2\)，末态 \(E'+E_e\)，其中 \(E_e=\sqrt{m^2c^4+p_e^2c^2}\)。三动量守恒 \(\mathbf p_e=\mathbf p_\gamma-\mathbf p'_\gamma\)，故 \(p_e^2c^2=E^2+E'^2-2EE'\cos\theta\)。
(2) 由 \(E_e=mc^2+E-E'\) 平方并代入动量式，约去公共项得 \(mc^2(E-E')=EE'(1-\cos\theta)\)。因此 \(E'=E/[1+(E/(mc^2))(1-\cos\theta)]\)，能量分母始终不小于一。
(3) 利用 \(E=hc/\lambda\)，得 \(\lambda'-\lambda=h(1-\cos\theta)/(mc)\)。电子动能 \(K_e=E-E'\)，在 \(\theta=\pi\) 最大，为 \(K_{\max}=E[1-1/(1+2E/(mc^2))]\)；正向散射 \(\theta=0\) 不转移能量。
(4) 若 \(E\ll mc^2\)，\(E'\simeq E[1-(E/(mc^2))(1-\cos\theta)]\)，反冲能相对光子能量很小，为Thomson低能近似。若 \(E\gg mc^2\) 且 \(\theta=\pi\)，\(E'\to mc^2/2\)，电子获得大部分初始光子能量。
